\scanpage{8}
\mspage{4}
% archive sheet number: 7

\uline{NB} ``Figure'' signifie ``figure finie'' -- i.e. on travaille dans l'atelier \inserted{$\unsure{\mathscr F|\mathcal F}\subset\unsure{\mathscr P|\mathcal P}(\unsure{\mathscr M|\mathcal M})$} des figures finies.

On trouve ainsi
\[
\begin{aligned}
\unsure{\operatorname{Pr\acute{e}fig}|\operatorname{Prefig}}(\mathscr M)&=\operatorname{Drap}^*(\mathscr M,\leq)\\
\operatorname{Fig}(\mathscr M)&=\operatorname{Drap}^*(\mathscr M,\leq)\cap\mathscr P_f(\mathscr M,\unsure{\unlhd|\trianglelefteq})
\end{aligned}
\]
Les figures sont les parties finies de $\mathscr M$, tot. ordonnées pour $\leq$, et fermées pour $\unlhd$.

Pour l'étude des supports constructibles et de la relation de complémentarité, il sera utile d'introduire, pour toute polyfigure, une ``polyfigure complémentaire. Quand $L$ n'a pas de plus petit \uline{et} de plus grand élément, ceci fait sortir des polyfigures proprement dites, en nous faisant intervenir des \deleted{\unsure{seg}} ``intervalles infinis'' $I_{<a}$, $I_{>a}$.

Même si $\deleted{\unsure{E}}L$ a plus petit et plus grand élément, l'introduction de $I_{<a}$, $I_{>a}$ rend le formalisme des polyfigures plus élégant, car la complémentaire $\unsure{\Phi'|\mathscr F'}$ d'une polyfigure $\Phi$ aura $\delta\Phi'\subset\delta\Phi$, alors que \deleted{en général} dans le contexte des polyfigures ``ordinaires'', la définition imposée par les circonstances donnerait seulement $\delta\Phi'\subset\delta\Phi\cup\underbrace{\{\omega_0,\omega_1\}}_{\unsure{\eta}}$, où $\omega_0,\omega_1$ sont le plus petit et plus grand élément. (De façon précise, on aurait $\eta=(\Phi\cap\eta)\amalg(\Phi'\cap\eta)$ -- tout point de $\eta$ (extr. de $L$) appartient à une des deux polyfigures
